\documentclass[10pt,a4paper]{amsart}
\usepackage[margin=19mm]{geometry}
\usepackage{amsmath,amssymb,mathtools}
\usepackage[scr=rsfs,cal=euler]{mathalfa}
\usepackage{xfrac}
\usepackage[unicode,hidelinks,bookmarks=false]{hyperref}
\newcommand{\ie}{i.e.\ }
\newcommand{\cf}{cf.\ }
\newcommand{\resp}{respectively\ }
\newcommand{\Subsection}{\subsection}
\providecommand{\nobreakdash}{\nobreak\hbox{-}}
\setlength{\emergencystretch}{2em}
\multlinegap0pt
\usepackage{amssymb}
\usepackage[shortlabels]{enumitem}
\usepackage{tikz}
\usepackage{mathtools}
\usepackage[normalem]{ulem}
\usepackage{xcolor}
\newtheorem{theorem}{Theorem}
\newcommand{\sPP}{\mathsf{sPP}}
\newcommand{\spp}{\sPP}
\title{Weight decomposition for toroidal abelian fibrations}
\author{Younghan Bae, Jeremy Feusi, Aitor Iribar L\'opez, Sam Molcho}
\date{Source: arXiv:2608.14337v1 (CC BY 4.0)}
\begin{document}
\maketitle

\noindent\emph{Source and licence.} Weight decomposition for toroidal abelian fibrations. Version arXiv:2608.14337v1 (CC BY 4.0), distributed by arXiv under the Creative Commons Attribution 4.0 International licence, http://creativecommons.org/licenses/by/4.0/, as recorded in the arXiv licence metadata for this exact version and shown on the abstract page https://arxiv.org/abs/2608.14337v1. Full licence notice: \texttt{licenses/arxiv-2608.14337-CC-BY-4.0.txt}.\par\smallskip

\noindent\textit{Editorial scope.} Counted unit: the article's theorem thm:polynomiality (source0.tex lines 2197--2203). The article's complete original proof of it is delivered with it (source0.tex lines 2959--3033, one \texttt{proof} environment of 75 lines, which the article places away from the statement). No mathematics is written, repaired or re-derived: the statement and the proof are the article's own lines, in the article's own notation, and no retained line is substituted or re-flowed.  No further source-local context is needed: every cross-reference of every retained line resolves inside the two delivered spans.  Nothing was omitted from any retained span, so the package carries no omission record.  No retained line contains a citation, so the wrapper carries no bibliography and no bibliography entry.  No retained line contains a figure environment, an image-inclusion command or drawing code, so no image is delivered and the package declares no asset dependency.  Editorial formatting only, in addition: an AMS \texttt{amsart} preamble that restates the article's own declarations and macros used by the retained lines, namely the environments \texttt{theorem} on separate counters, together with the standard packages those lines need.  The retained environments print in this document as Theorem 1 in order of appearance, whereas the article prints the same items with the numbering of its own full document.  The complete native source of the article as distributed in the arXiv e-print for this exact version is bundled unchanged as \texttt{sources/207660/source0.tex}.
\begin{theorem}
	\label{thm:polynomiality}
	Let $f\in \spp^{\ast}(\tilde{\Sigma}^{(s)})[m]^{\mathsf{gl}}$ and set $d:=\deg_{m,\mathsf{x}_s}(f)$. For every maximal cone
	$\sigma$ of $\tilde{\Sigma}^{(s)}$, the restriction
	$$\left.\mathfrak N_{N,m}^\ast(f)\right|_\sigma$$
	is a polynomial in $N$ of degree at most $d$.
\end{theorem}

\begin{proof}[Proof of Theorem \ref{thm:polynomiality}]
	Fix $\rho\in \mathfrak S_{s-1}$ and $0\le r\le s-1$. For simplicity, set
	\[
		c(m):=f|_{\sigma_{\rho_0}}(m)(\mathsf{x}_s\to 0).
	\]
	The continuity condition in the definition of $\spp^{\ast}(\tilde{\Sigma}^{(s)})[m]^{\mathsf{gl}}$ gives
	\[
		f|_{\sigma_{\rho_{s-1}}}(m)(\mathsf{x}_s\to \mathsf{t})=c(m+1).
	\]
	Thus we may write
	\[
		f|_{\sigma_{\rho_0}}(m)=c(m)+\mathsf{x}_sR_0(m),\qquad
		f|_{\sigma_{\rho_{s-1}}}(m)=c(m+1)+(\mathsf{x}_s-\mathsf{t})R_s(m)
	\]
	with $R_0$ and $R_s$ in $\mathbb{Q}[m,\mathsf{x}_1,\dots,\mathsf{x}_s,\mathsf{t}]$. Substituting this into $S_{i,r}$, we rewrite
    {\allowdisplaybreaks
	\begin{align*}
		S_{i,r}={}&\left(1+\frac{(i-1)(\mathsf{t}-\mathsf{x}_s)}{\mathsf{z}_i^{(N)}}\right)c(Nm+i-1)
		+\frac{(N-i)\mathsf{x}_s}{\mathsf{t}-\mathsf{z}_i^{(N)}}c(Nm+i)\\
				  &+(i-1)(\mathsf{t}-\mathsf{x}_s)R_0(m\to Nm+i-1,\mathsf{x}_s\to \mathsf{z}_i^{(N)})
				  +\sum_{h=1}^{r}(i-1)(\mathsf{t}-\mathsf{x}_s)H_{\rho,h}^f(m\to Nm+i-1,\mathsf{x}_s\to \mathsf{z}_i^{(N)})\notag\\
				  &+f|_{\sigma_{\rho_r}}(m\to Nm+i-1,\mathsf{x}_s\to \mathsf{z}_i^{(N)})-c(Nm+i-1)
				  +\sum_{h=r+1}^{s-1}(N-i)\mathsf{x}_sH_{\rho,h}^f(m\to Nm+i-1,\mathsf{x}_s\to \mathsf{z}_i^{(N)})\notag\\
				  &-(N-i)\mathsf{x}_sR_s(m\to Nm+i-1,\mathsf{x}_s\to \mathsf{z}_i^{(N)}).
	\end{align*}}
	The coefficient of $c(Nm+i-1)$ simplifies as
	\[
		1+\frac{(i-1)(\mathsf{t}-\mathsf{x}_s)}{\mathsf{z}_i^{(N)}}=\frac{(N-i+1)\mathsf{x}_s}{\mathsf{z}_i^{(N)}}.
	\]
	Using $\mathsf{t}-\mathsf{z}_i^{(N)}=-\mathsf{z}_{i+1}^{(N)}$, the sum of the terms involving
	$c$ is
	\[
		\sum_{i=1}^N\left(\frac{(N-i+1)\mathsf{x}_s}{\mathsf{z}_i^{(N)}}c(Nm+i-1)
		-\frac{(N-i)\mathsf{x}_s}{\mathsf{z}_{i+1}^{(N)}}c(Nm+i)\right),
	\]
	which telescopes to
	\[
		c(Nm).
	\]
	All remaining terms are polynomial in $(i,N)$. Hence polynomiality in $N$ follows.

	It remains to prove the degree bound. If $d=0$, all $H_{\rho,h}^{f}$ and the functions $R_0,R_s$ vanish, $c$ is independent
	of $m$ and
	\[
		\left.\mathfrak N_{N,m}^\ast(f)\right|_{\sigma_{\rho_r}}=c(Nm),
	\]
	which has degree at most $0$ in $N$.

	Assume $d\ge 1$. Note that the degree of $c$ in $m$ is at most $d$ and hence the contribution $c(Nm)$ has
	degree at most $d$ in $N$. The remaining terms in the above expression for $S_{i,r}$ are clearly polynomial in $i$ and $N$ of total
	degree at most $d$. Hence it suffices to show that the degree $d$ part of the remainder of $S_{i,r}$ vanishes.

	For $0\le h\le s-1$, let $G_h(m,\mathsf{x}_s)$ be the homogeneous degree-$d$ part in $m$ and $\mathsf{x}_s$ of $f|_{\sigma_{\rho_h}}(m)$, and let $C(m)=G_0(m,0)$. The homogeneous degree-$(d-1)$ parts in $m$ and $\mathsf{x}_s$ of $R_0$, $R_s$, and $H_{\rho,h}^f$ are respectively
	\[
		\frac{G_0-C}{\mathsf{x}_s},\qquad \frac{G_{s-1}-C}{\mathsf{x}_s},\qquad \frac{G_h-G_{h-1}}{\mathsf{x}_s}.
	\]
	Since the substitutions $m\to Nm+i-1$ and $\mathsf{x}_s\to \mathsf{z}_i^{(N)}$ are of total degree 1 in $i$ and $N$, it follows that the
	possible degree-$d$ contribution in $i$ and $N$ to $S_{i,r}$ (omitting the $c$ terms as explained above) is
	the degree $d$ part of
	\[
		\frac{(i-1)(\mathsf{t}-\mathsf{x}_s)}{\mathsf{z}_i^{(N)}}(G_0-C)
		+\frac{(i-1)(\mathsf{t}-\mathsf{x}_s)}{\mathsf{z}_i^{(N)}}\sum_{h=1}^r(G_h-G_{h-1})
		+(G_r-C)
	\]
	\[
		+\frac{(N-i)\mathsf{x}_s}{\mathsf{z}_i^{(N)}}\sum_{h=r+1}^{s-1}(G_h-G_{h-1})
		-\frac{(N-i)\mathsf{x}_s}{\mathsf{z}_i^{(N)}}(G_{s-1}-C).
	\]
	The sums telescope, so this equals
	\[
		(G_r-C)\left(1+\frac{(i-1)(\mathsf{t}-\mathsf{x}_s)}{\mathsf{z}_i^{(N)}}-\frac{(N-i)\mathsf{x}_s}{\mathsf{z}_i^{(N)}}\right)
		=\frac{\mathsf{x}_s(G_r-C)}{\mathsf{z}_i^{(N)}}.
	\]
	This has total degree $d-1$ in $(m,\mathsf{x}_s,i,N)$ as required. Substituting $m$ and $\mathsf{x}_s$ as above and summing over $i=1,\dots,N$ gives degree at most $d$ in $N$.
\end{proof}

\end{document}
